% CREATED BY Kyriaki Antoniadou-Plytaria, 2021
% MODIFIED BY Igor S Peretta, 2023

\usepackage{xspace}
\newcommand{\ppcversao}{2026-1}
\newcommand{\nomecursonovo}{Engenharia de Computação com Inteligência Artificial Aplicada\xspace}%
\newcommand{\fraseaplicada}{O termo ``Aplicada'' na denominação reforça a vocação do curso para formar engenheiros capazes de aplicar técnicas de IA na resolução de problemas reais de engenharia, tais como cibersegurança, sistemas embarcados inteligentes, otimização de redes, processamento de sinais e sistemas de controle adaptativos.\xspace}%

% BASIC SETTINGS
\usepackage{moreverb}								% List settings
\usepackage{textcomp}								% Fonts, symbols etc.
\usepackage{lmodern}								% Latin modern font
\usepackage{helvet}									% Enables font switching
\usepackage[T1]{fontenc}							% Output settings
\usepackage[brazil]{babel}							% Language settings
\usepackage[utf8]{inputenc}							% Input settings
\usepackage{amsmath}								% Mathematical expressions (American mathematical society)
\usepackage{amssymb}								% Mathematical symbols (American mathematical society)
\usepackage{graphicx}								% Figures
\usepackage{subfig}									% Enables subfigures
\numberwithin{equation}{chapter}					% Numbering order for equations
\numberwithin{figure}{chapter}						% Numbering order for figures
\numberwithin{table}{chapter}						% Numbering order for tables
\usepackage{listings}								% Enables source code listings
% \usepackage{chemfig}								% Chemical structures
\usepackage{multirow}
\usepackage[table]{xcolor}
\usepackage[top=3cm, bottom=3cm,
			inner=3cm, outer=3cm]{geometry}			% Page margin lengths			
\usepackage{eso-pic}								% Create cover page background
\newcommand{\backgroundpic}[3]{
	\put(#1,#2){
	\parbox[b][\paperheight]{\paperwidth}{
	\centering
	\includegraphics[width=\paperwidth,height=\paperheight,keepaspectratio]{#3}}}}
\usepackage{float} 									% Enables object position enforcement using [H]
\usepackage{parskip}								% Enables vertical spaces correctly 
\usepackage{wrapfig}
\usepackage{hyperref}
\usepackage{tikz}
\usetikzlibrary{fit}
\usepackage{enumerate}
\usepackage{csquotes}
\usepackage{soul}
\usepackage{caption}
% \usepackage{afterpage}
\usepackage{makecell}


% DEFINE THESIS TYPE
% \def\ThesisType{B} % M - master, B - bachelor

% DEFINE INSTITUTION LOCATION
% \def\InstitutionLocation{C} % C - Chalmers only, G - Chalmers and GU combined

% COLOR FOR HEADERS
% COLOR FOR HEADERS
\definecolor{AzulClaro}{RGB}{41,102,189}
\definecolor{AzulEscuro}{RGB}{14,64,151}
\definecolor{CinzaClaro}{RGB}{232,231,231}
\definecolor{CinzaMeio}{RGB}{172,175,180}
\definecolor{CinzaEscuro}{RGB}{113,119,130}
\definecolor{Verde}{RGB}{1,66,10}
\definecolor{Amarelo}{RGB}{232,171,2}
\definecolor{Laranja}{RGB}{204,85,0}
\definecolor{AzulMarca}{RGB}{220,230,241}
\definecolor{AmareloVivo}{RGB}{255,255,0}


\hypersetup{
    colorlinks=true,
    linkcolor=AzulClaro,
    filecolor=AzulClaro,
    citecolor = AzulEscuro,
    urlcolor=AzulClaro,
}
% OPTIONAL SETTINGS (DELETE OR COMMENT TO SUPRESS)

\usepackage[explicit]{titlesec}
%%% CHAPTER
\titleformat{\chapter}[display]
  {\normalfont}
  {}
  {0pt}
  {%
    \begin{tikzpicture}
    \node[
      fill=AzulEscuro,
      minimum width=1.1cm,
      minimum height=1.1cm,
      rounded corners,
      anchor=west,
      font=\color{white}\fontsize{18}{24}\selectfont\bfseries
      ]
      % at ([xshift=6pt]chapname.south)
      (chapnumber)
      {\thechapter};
    \node[
      anchor=west,
      text width=\textwidth-4cm,
      font=\color{AzulEscuro}\bfseries\Large
      ] 
      at ([xshift=10pt]chapnumber.east)
      {#1};
    \fill[
      overlay,
      draw=none,
      line width=0pt,
      rounded corners=0pt,
      left color=AzulEscuro,
      right color=AzulClaro!25
      ]
      ([yshift=12pt]chapnumber.north west) rectangle ++(\textwidth,-2pt);  
    \end{tikzpicture}%
  }
\titleformat{name=\chapter,numberless}[display]
  {\normalfont}
  {}
  {0pt}
  {%
    \begin{tikzpicture}
    \node[
      anchor=west,
      inner sep=0pt,
      outer sep=0pt,
      text width=\textwidth,
      font=\color{AzulEscuro}\bfseries\Large
      ]
      (chaptitle) 
      {#1};
    \fill[
      overlay,
      draw=none,
      line width=0pt,
      rounded corners=1pt,
      left color=AzulEscuro,
      right color=AzulClaro!25
      ]
      ([yshift=-15pt]chaptitle.south west) rectangle ++(\textwidth,-2pt);  
    \end{tikzpicture}%
  }
%%% SECTION
\titleformat{\section}[display]
  {\normalfont}
  {}
  {0pt}
  {%
    \begin{tikzpicture}
    \node[
      fill=AzulClaro,
      minimum width=0.9cm,
      minimum height=0.9cm,
      rounded corners,
      anchor=west,
      font=\color{white}\fontsize{12}{14}\selectfont\bfseries
      ]
      % at ([xshift=4pt]square.south)
      (secnumber)
      {\thesection};
    \node[
      anchor=west,
      text width=\textwidth-4cm,
      font=\color{AzulClaro}\bfseries\large
      ] 
      at ([xshift=8pt]secnumber.east)
      {#1};
    \end{tikzpicture}%
  }
\titleformat{name=\section,numberless}[display]
  {\normalfont}
  {}
  {0pt}
  {%
    \begin{tikzpicture}
    \node[
      anchor=west,
      inner sep=0pt,
      outer sep=0pt,
      text width=\textwidth,
      font=\color{AzulClaro}\bfseries\normalsize
      ]
      (sectitle) 
      {#1};
    \end{tikzpicture}%
  }
% \titlespacing*{\section}{0pt}{0.5ex plus 1ex minus .2ex}{2.3ex plus .2ex}
% Disable automatic indentation (equal to using \noindent)
% \setlength{\parindent}{0cm}                         

\usepackage{tabularray}
% \usepackage{array}
% \let\oldtabular\tabular
% \let\endoldtabular\endtabular
% \renewenvironment{tabular}{\rowcolors{2}{grey!25}{white}\oldtabular}{\endoldtabular}
% \newcommand{\tabh}[2]{\textcolor{white}{\multicolumn{1}{>{\centering\arraybackslash}p{#1}}{\textbf{#2}}}}
\NewTblrTheme{ecp}{%
    \DefTblrTemplate{caption-tag}{normal}{\textbf{\textcolor{AzulEscuro}{Quadro\hspace{0.25em}\thetable}}}
    \DefTblrTemplate{caption-sep}{normal}{\textbf{\textcolor{AzulEscuro}{:}}\enskip}
    \DefTblrTemplate{caption-text}{normal}{\InsertTblrText{caption}}
    \DefTblrTemplate{contfoot-text}{normal}{\tiny \textcolor{AzulClaro}{Continua na próxima página}}
    \DefTblrTemplate{conthead-text}{normal}{\tiny \textcolor{AzulClaro}{(Continuação)}}
    \SetTblrTemplate{caption-tag}{normal}
    \SetTblrTemplate{caption-sep}{normal}
    \SetTblrTemplate{caption-text}{normal}
    \SetTblrTemplate{contfoot-text}{normal}
    \SetTblrTemplate{conthead-text}{normal}
}

% Caption settings (aligned left with bold name)
\usepackage[labelfont=bf, textfont=normal,
			justification=justified,
			singlelinecheck=false,
            % tablename=Quadro
            ]{caption} 		
% \captionsetup[table]{name=Quadro}
% \renewcommand{\tableautorefname}{Quadro}
		  	
% Activate clickable links in table of contents  	
% \usepackage{hyperref}								
% \hypersetup{colorlinks, citecolor=black,
%    		 	filecolor=black, linkcolor=black,
%     		urlcolor=black}


% Define the number of section levels to be included in the t.o.c. and numbered	(3 is default)	
\setcounter{tocdepth}{5}							
\setcounter{secnumdepth}{5}	


% % Chapter title settings
% \usepackage{titlesec}		
% \titleformat{\chapter}[display]
%   {\Huge\bfseries\filcenter}
%   {{\fontsize{50pt}{1em}\vspace{-4.2ex}\selectfont \textnormal{\thechapter}}}{1ex}{}[]


% Header and footer settings (Select TWOSIDE or ONESIDE layout below)
\usepackage{fancyhdr}								
\pagestyle{fancy}  
\renewcommand{\chaptermark}[1]{\markboth{\thechapter.\space#1}{}} 


% Select one-sided (1) or two-sided (2) page numbering
\def\layout{2}	% Choose 1 for one-sided or 2 for two-sided layout
% Conditional expression based on the layout choice
\ifnum\layout=2	% Two-sided
    \fancyhf{}			 						
	\fancyhead[LE,RO]{\nouppercase{ \leftmark}}
	\fancyfoot[LE,RO]{\thepage}
	\fancypagestyle{plain}{			% Redefine the plain page style
	\fancyhf{}
	\renewcommand{\headrulewidth}{0pt} 		
	\fancyfoot[LE,RO]{\thepage}}	
\else			% One-sided  	
  	\fancyhf{}					
	\fancyhead[C]{\nouppercase{ \leftmark}}
	\fancyfoot[C]{\thepage}
\fi


% % Enable To-do notes
% \usepackage[textsize=tiny]{todonotes}   % Include the option "disable" to hide all notes
% \setlength{\marginparwidth}{2.5cm} 

% \newcommand{\todoreview}[1]{ \textcolor{Laranja}{\ul{#1}}\todo{Rever!}}

% Supress warning from Texmaker about headheight
\setlength{\headheight}{15pt}		


% Referência para autor do template original:
% Tag command (invisible)
% \newcommand{\tagtemp}{\textcolor{white}{, template by Kyriaki Antoniadou-Plytaria}}

% Allow a table to split
% \usepackage{longtable}

\usepackage{footnotehyper}
\makesavenoteenv{longtblr}
\usepackage[inline]{enumitem}

\newcommand{\listspace}{\setlist[enumerate]{itemsep=5pt,parsep=0pt,topsep=0pt,partopsep=0pt}}

\usepackage{lscape}
\usepackage{xifthen}

\newcommand\nohyph{\hyphenpenalty=10000\relax\exhyphenpenalty=10000\relax} 

\usepackage{etoolbox}
\AtBeginDocument{%
    \renewcommand{\familydefault}{\sfdefault}
    \normalfont
}


